% !TeX root = ../main.tex
\section{Đánh giá kết quả}\label{sec:danhgia}

\subsection{Kho dữ liệu}

Bảng~\ref{tab:danhgiakho} tổng hợp các phép kiểm tra đã chạy trên hệ thống thật ở phần hiện thực, cùng kết quả đo được.

\begingroup\small\setstretch{1.1}
\begin{longtable}{|P{3.0cm}|>{\raggedright\arraybackslash}p{5.2cm}|>{\raggedright\arraybackslash}p{5.5cm}|}
\caption{Kết quả kiểm tra kho dữ liệu}\label{tab:danhgiakho}\\
\hline
\textbf{Tiêu chí} & \textbf{Cách kiểm tra} & \textbf{Kết quả}\\ \hline \endfirsthead
\hline
\textbf{Tiêu chí} & \textbf{Cách kiểm tra} & \textbf{Kết quả}\\ \hline \endhead
Thời gian ETL & Nhật ký \code{meta.etl\_run\_log} của một lần nạp đầy đủ & 8,7 giây cho 1,55 triệu dòng nguồn (Bảng~\ref{tab:etlbuoc})\\ \hline
Khớp số liệu với nguồn & 20 tiêu chí đối soát trong \code{40\_dq\_checks.sql} & 20/20 tiêu chí đạt\\ \hline
Toàn vẹn tham chiếu & Ràng buộc khóa ngoại và bản ghi ``không xác định'' ($-1$) ở mọi bảng chiều & Mọi dòng của \code{fact\_sales} trỏ đến một bản ghi chiều hợp lệ\\ \hline
Truy vết bản ghi bị loại, bị sửa & Bảng \code{meta.rejected\_record} & 775 đơn bị loại và 1.458 bản ghi được sửa, mỗi bản ghi kèm lý do\\ \hline
Chạy lại nhiều lần & So số dòng và tổng tiền giữa các lần chạy & Không đổi, không trùng dữ liệu\\ \hline
Lỗi giữa chừng & \code{verify\_rollback.py} cố ý gây lỗi trước khi xác nhận giao dịch & Data Mart trở lại đủ 112.650 dòng\\ \hline
Tốc độ truy vấn OLAP & \code{EXPLAIN ANALYZE}, trung vị 5 lần chạy & Chậm nhất 0,72 giây (CUBE vùng và năm, Bảng~\ref{tab:thoigiantruyvan})\\ \hline
Phân quyền & Role \code{analyst\_ro}; thông tin kết nối trong \code{.env} & Người dùng phân tích chỉ có quyền \code{SELECT}\\ \hline
Metadata & Schema \code{meta} & Có từ điển dữ liệu và nguồn gốc của từng cột\\ \hline
Đúng số liệu trên dashboard & Đối chiếu tám chỉ số trên dashboard với kết quả OLAP và DSS & Cả tám chỉ số khớp (Bảng~\ref{tab:doichieu})\\ \hline
Thời gian phản hồi dashboard & Mở từng trang trong tiến trình mới bằng AppTest của Streamlit & Lần đầu 0,4-4,7 giây, lần sau dưới 0,1 giây; đạt yêu cầu tối đa 5 giây\\ \hline
Tính khả dụng & Các thao tác lọc, drill-down, pivot, what-if trên dashboard & Người dùng chọn bằng thanh trượt, ô chọn và nút bấm, không cần viết SQL; tên chỉ số và chiều bằng tiếng Việt, số theo định dạng Việt Nam\\ \hline
Phân quyền dashboard & Thử \code{UPDATE}, \code{DELETE} bằng tài khoản \code{dashboard\_app} & Bị từ chối; dashboard chỉ được đọc và ghi thêm vào \code{dss.decision\_log}\\ \hline
\end{longtable}
\endgroup

Mọi nhu cầu theo dõi của bốn nhóm người dùng ở Bảng~\ref{tab:quyetdinh} đều có ít nhất một truy vấn OLAP trả lời: doanh thu theo thời gian, khu vực và ngành hàng (roll-up, slice, dice); tỷ lệ trễ theo bang và tuyến cùng thời gian từng công đoạn (drill-down, ma trận tuyến); điểm đánh giá theo ngành hàng, vùng và tình trạng giao hàng (pivot); người bán doanh thu cao nhưng chất lượng thấp và phân bố người bán so với khách hàng. Dashboard ở mục~\ref{sec:dashboard} đưa các chỉ số này lên giao diện để người dùng tự lọc và xem.

\subsection{Hệ hỗ trợ quyết định}

DSS đạt cả bốn chỉ tiêu đã đặt (Bảng~\ref{tab:datchitieu}): khi cảnh báo 10\% số đơn bàn giao mỗi ngày, mô hình bắt được 41,7\% số đơn sẽ trễ, và 44,5\% khi kết hợp với luật ``đã quá hạn''; kết quả ổn định qua cả ba tháng test, và đơn mức ĐỎ trễ nhiều gấp 11,7 lần đơn mức XANH. Đợt tìm kiếm mở rộng (124 cấu hình, thêm đặc trưng tồn đọng và mô hình kết hợp) cho kết quả test trong khoảng 39,6-42,6\%, nằm trong sai số của mô hình được chọn, cho thấy kết quả này ổn định chứ không phụ thuộc vào một bộ tham số may mắn. Phân tích chi phí cho người quản lý một quy tắc dùng được ngay: hành động rẻ như báo trước cho khách nên làm cho mọi đơn ĐỎ, còn hành động đắt như tặng voucher chỉ nên dành cho các đơn vượt ngưỡng $p^{*}$. Với giả định chi phí mặc định, chính sách này có lợi ích ròng dương (688 BRL trên ba tháng test), trong khi tặng voucher cho cả 15\% đơn rủi ro nhất thì lỗ.
